\chapter{Annexes}

\section{Annexe A : Code source du firmware ESP32}

\subsection{Programme principal}

\begin{lstlisting}[language=C++,
caption={Firmware Andon ESP32 - complet},
label={lst:firmware_complet}]
#include <WiFi.h>
#include <PubSubClient.h>
#include <ArduinoJson.h>

const char* SSID = "SEWS_IOT";
const char* PASSWORD = "****";
const char* MQTT_SERVER = "192.168.1.100";
const int MQTT_PORT = 1883;
const int LIGNE_ID = 1;
const int BOUTON_PIN = 4;
const int LED_VERTE = 16;
const int LED_ROUGE = 17;
const unsigned long DEBOUNCE = 200;

WiFiClient wifiClient;
PubSubClient mqtt(wifiClient);
unsigned long dernierAppui = 0;

void connectWiFi() {
  WiFi.begin(SSID, PASSWORD);
  while (WiFi.status() != WL_CONNECTED) {
    delay(500);
  }
}

void connectMQTT() {
  while (!mqtt.connected()) {
    if (mqtt.connect("esp32_ligne1")) {
      String topic = "sews/command/" 
        + String(LIGNE_ID) + "/led";
      mqtt.subscribe(topic.c_str());
    } else {
      delay(5000);
    }
  }
}

void envoyerAlerte() {
  StaticJsonDocument<256> doc;
  doc["ligneId"] = LIGNE_ID;
  doc["type"] = "panne";
  doc["urgence"] = "haute";
  char buffer[256];
  serializeJson(doc, buffer);
  String topic = "sews/andon/" 
    + String(LIGNE_ID) + "/alert";
  mqtt.publish(topic.c_str(), buffer, true);
}

void callbackMQTT(char* topic, byte* payload,
    unsigned int len) {
  StaticJsonDocument<128> doc;
  deserializeJson(doc, payload, len);
  digitalWrite(LED_ROUGE, 
    doc["rouge"] ? HIGH : LOW);
  digitalWrite(LED_VERTE, 
    doc["verte"] ? HIGH : LOW);
}

void setup() {
  Serial.begin(115200);
  pinMode(BOUTON_PIN, INPUT);
  pinMode(LED_VERTE, OUTPUT);
  pinMode(LED_ROUGE, OUTPUT);
  digitalWrite(LED_VERTE, HIGH);
  connectWiFi();
  mqtt.setServer(MQTT_SERVER, MQTT_PORT);
  mqtt.setCallback(callbackMQTT);
  connectMQTT();
}

void loop() {
  if (!mqtt.connected()) connectMQTT();
  mqtt.loop();
  if (digitalRead(BOUTON_PIN) == HIGH) {
    unsigned long now = millis();
    if (now - dernierAppui > DEBOUNCE) {
      dernierAppui = now;
      envoyerAlerte();
      digitalWrite(LED_VERTE, LOW);
      digitalWrite(LED_ROUGE, HIGH);
    }
  }
}
\end{lstlisting}

\section{Annexe B : Script SQL}

\begin{lstlisting}[language=SQL,
caption={Initialisation de la base},
label={lst:sql_init}]
CREATE TABLE lignes (
    id SERIAL PRIMARY KEY,
    nom VARCHAR(100) NOT NULL,
    numero INTEGER UNIQUE NOT NULL,
    statut VARCHAR(20) DEFAULT 'active'
);

CREATE TABLE evenements (
    id SERIAL PRIMARY KEY,
    ligne_id INTEGER REFERENCES lignes(id),
    type VARCHAR(50) NOT NULL,
    statut VARCHAR(20) DEFAULT 'en_cours',
    description TEXT,
    date_debut TIMESTAMP DEFAULT NOW(),
    date_fin TIMESTAMP,
    technicien_id INTEGER 
      REFERENCES utilisateurs(id),
    compte_rendu TEXT
);

CREATE TABLE utilisateurs (
    id SERIAL PRIMARY KEY,
    nom VARCHAR(100) NOT NULL,
    prenom VARCHAR(100) NOT NULL,
    email VARCHAR(255) UNIQUE NOT NULL,
    role VARCHAR(50) NOT NULL,
    mot_de_passe VARCHAR(255) NOT NULL
);

INSERT INTO lignes (nom, numero) VALUES
  ('Ligne T1', 1),
  ('Ligne T2', 2),
  ('Ligne T3', 3),
  ('Ligne T4', 4),
  ('Ligne T5', 5),
  ('Ligne T6', 6),
  ('Ligne T7', 7),
  ('Ligne T8', 8);
\end{lstlisting}

\section{Annexe C : package.json}

\begin{lstlisting}[caption={Dependances backend},
label={lst:package_backend}]
{
  "name": "sews-andon-server",
  "dependencies": {
    "express": "^4.18.2",
    "socket.io": "^4.7.2",
    "mqtt": "^5.3.0",
    "pg": "^8.11.3",
    "jsonwebtoken": "^9.0.2",
    "bcrypt": "^5.1.1"
  }
}
\end{lstlisting}

\section{Annexe D : Pinout ESP32}

\begin{table}[htbp]
\centering
\begin{tabular}{l l l}
\toprule
\textbf{Broche} & \textbf{GPIO} & \textbf{Fonction} \\
\midrule
D4 & GPIO 4 & Bouton Andon \\
D16 & GPIO 16 & LED verte \\
D17 & GPIO 17 & LED rouge \\
VIN & 5V & Alimentation \\
GND & GND & Masse \\
\bottomrule
\end{tabular}
\caption{Pinout ESP32}
\label{tab:pinout}
\end{table}